% !PRINTED 119
\weberpage{119}

Dies ist ein Mangel der \gesperrt{Cardanischen Formel}, die hiernach für sich noch nicht genügt, um die Wurzeln der cubischen Gleichung eindeutig zu geben. Diesem Mangel hat \gesperrt{Cayley} durch die folgende Darstellung abgeholfen. Man definire zwei neue Grössen $\xi,\eta$ durch die Gleichungen
\begin{equation}
u=\xi^2\eta,
\qquad
v=\xi\eta^2,
\tag{2}
\end{equation}
woraus sich durch Multiplication mit Benutzung von (1)
\begin{equation}
\xi\eta=\sqrt[3]{-\frac a3}
\tag{3}
\end{equation}
ergiebt; folglich ist, wenn dieser Werth in (2) eingesetzt und für $u,v$ ihre Ausdrücke \S.\,35, (8) substituirt werden,
\begin{equation}
\xi=\sqrt[3]{\frac{3b}{2a}+\sqrt{\frac{9b^2}{4a^2}+\frac a3}},\qquad
\eta=\sqrt[3]{\frac{3b}{2a}-\sqrt{\frac{9b^2}{4a^2}+\frac a3}},
\tag{4}
\end{equation}
und wenn man also jetzt die Wurzel $y$ der cubischen Gleichung in die Form
\begin{equation}
y=\xi\eta(\xi+\eta)
\tag{5}
\end{equation}
setzt, so erhält man einen Ausdruck, der, wenn man, von einander unabhängig, $\xi$ durch $\xi,\varepsilon\xi,\varepsilon^2\xi$, und $\eta$ durch $\eta,\varepsilon\eta,\varepsilon^2\eta$ ersetzt, nicht neun sondern \gesperrt{nur die drei Werthe}
\[
\begin{array}{c}
\xi^2\eta+\xi\eta^2,\\
\varepsilon\xi^2\eta+\varepsilon^2\xi\eta^2,\\
\varepsilon^2\xi^2\eta+\varepsilon\xi\eta^2
\end{array}
\]
annimmt.
\footnote{\gesperrt{Cayley}, Phil. Mag. vol. XXI, 1861. Collectet mathematical papers vol. V, Nr. 310.}% [sic] Weber druckt "Collectet" (Druckfehler fuer "Collected"); quellentreu reproduziert.



\begin{center}
{\large \S.\ 37.}

\vspace{0.5em}
{\large\gesperrt{Die biquadratische Gleichung.}}
\end{center}

Ein ähnlicher Weg, wie bei der Lösung der cubischen Gleichung durch die \gesperrt{Cardanische Formel}, lässt sich zur Lösung der Gleichung vierten Grades
\begin{equation}
 y^4+ay^2+by+c=0
\tag{1}
\end{equation}
einschlagen. Wir setzen
%CONT>>
